% GENERATED FILE. Edit canonical lessons, then run npm run generate:books.
% canonical-source-hash: fnv1a64:38f5b864de896376
% canonical-lessons: PA-C348-saz, PA-C348-tabla, PA-C348-dhun, PA-C348-riaz, PA-C348-gunjna

\chapter{Instruments I Like}
\label{ch:pa-a2-instruments-i-like}
\hlchaptermodality{348}

\begin{chapteropening}
\textbf{By the end of this chapter:} \emph{I can name the instruments I like and say how I practise.}

Complete the last lesson of chapter 348: I can name the instruments I like and say how I practise.
\end{chapteropening}

\section[\texorpdfstring{\pa{ਸਾਜ਼} (sāz)}{ਸਾਜ਼ (sāz)}]{\pa{ਸਾਜ਼} (sāz) — a musical instrument}

\label{lesson:PA-C348-saz}



\begin{quote}
Before the new one: say the Punjabi for a drumbeat, then the Punjabi for a clap.
\end{quote}

\subsection*{You'll want to know: \pa{ਸਾਜ਼}}
\textbf{\pa{ਸਾਜ਼}} — \emph{sāz} — \textquotedblleft{}a musical instrument\textquotedblright{}.

\begin{quote}
\pa{ਮੈਨੂੰ} \pa{ਹਰ} \pa{ਸਾਜ਼} \pa{ਪਸੰਦ} \pa{ਹੈ।}
\end{quote}

\begin{quote}
\emph{mainū̃ har sāz pasand hai.}
\end{quote}

\begin{quote}
I like every musical instrument.
\end{quote}

A masculine noun.

\begin{cousinweb}
\textbf{\pa{ਸਾਜ਼}} (\emph{sāz}) is a loan, usually traced to Persian \emph{sāz}, \textquotedblleft{}instrument\textquotedblright{}.
\end{cousinweb}

\begin{grammarlens}[title={What you've built}]
One more word for this chapter.
\end{grammarlens}

\subsection*{Your turn}
\begin{itemize}
\raggedright
  \item \emph{Say it:} \emph{sāz}
  \item \emph{Say it:} \emph{sāz}, once more
  \item \emph{Say it:} \emph{mainū̃ har sāz pasand hai.} — I like every musical instrument.
  \item \emph{Recall:} say the Punjabi for a drumbeat, then the Punjabi for a clap, then say \emph{sāz} again
\end{itemize}

\begin{tcolorbox}[breakable,colback=teal!4,colframe=teal!35!black,title={Before you move on}]
What does \pa{ਸਾਜ਼} mean? (\textquotedblleft{}A musical instrument\textquotedblright{}.) Say it once more.
\end{tcolorbox}
\section[\texorpdfstring{\pa{ਤਬਲਾ} (tablā)}{ਤਬਲਾ (tablā)}]{\pa{ਤਬਲਾ} (tablā) — a tabla; a pair of hand drums}

\label{lesson:PA-C348-tabla}



\begin{quote}
Before the new one: say the Punjabi for a clap, then the Punjabi for a musical instrument.
\end{quote}

\subsection*{You'll want to know: \pa{ਤਬਲਾ}}
\textbf{\pa{ਤਬਲਾ}} — \emph{tablā} — \textquotedblleft{}a tabla; a pair of hand drums\textquotedblright{}.

\begin{quote}
\pa{ਅਮਨ} \pa{ਹਰ} \pa{ਰੋਜ਼} \pa{ਤਬਲਾ} \pa{ਵਜਾਉਂਦਾ} \pa{ਹੈ।}
\end{quote}

\begin{quote}
\emph{Aman har roz tablā vajāundā hai.}
\end{quote}

\begin{quote}
Aman plays the tabla every day.
\end{quote}

A masculine noun.

\begin{cousinweb}
\textbf{\pa{ਤਬਲਾ}} (\emph{tablā}) is a loan, usually traced through Persian to Arabic \emph{ṭabl}, \textquotedblleft{}a drum\textquotedblright{}.
\end{cousinweb}

\begin{grammarlens}[title={What you've built}]
One more word for this chapter.
\end{grammarlens}

\subsection*{Your turn}
\begin{itemize}
\raggedright
  \item \emph{Say it:} \emph{tablā}
  \item \emph{Say it:} \emph{tablā}, once more
  \item \emph{Say it:} \emph{Aman har roz tablā vajāundā hai.} — Aman plays the tabla every day.
  \item \emph{Recall:} say the Punjabi for a clap, then the Punjabi for a musical instrument, then say \emph{tablā} again
  \item \emph{Write it:} copy \textbf{\pa{ਤਬਲਾ}} once, checking the whole word
\end{itemize}

\begin{tcolorbox}[breakable,colback=teal!4,colframe=teal!35!black,title={Before you move on}]
What does \pa{ਤਬਲਾ} mean? (\textquotedblleft{}A tabla; a pair of hand drums\textquotedblright{}.) Say it once more.
\end{tcolorbox}
\section[\texorpdfstring{\pa{ਧੁਨ} (dhun)}{ਧੁਨ (dhun)}]{\pa{ਧੁਨ} (dhun) — a tune; a melody}

\label{lesson:PA-C348-dhun}



\begin{quote}
Before the new one: say the Punjabi for a musical instrument, then the Punjabi for a tabla.
\end{quote}

\subsection*{You'll want to know: \pa{ਧੁਨ}}
\textbf{\pa{ਧੁਨ}} — \emph{dhun} — \textquotedblleft{}a tune; a melody\textquotedblright{}.

\begin{quote}
\pa{ਉਸ} \pa{ਗੀਤ} \pa{ਦੀ} \pa{ਧੁਨ} \pa{ਬਹੁਤ} \pa{ਮਿੱਠੀ} \pa{ਹੈ।}
\end{quote}

\begin{quote}
\emph{us gīt dī dhun bahut miṭṭhī hai.}
\end{quote}

\begin{quote}
The tune of that song is very sweet.
\end{quote}

A feminine noun.

\begin{cousinweb}
\textbf{\pa{ਧੁਨ}} (\emph{dhun}) is inherited from Sanskrit \emph{dhuni}, \textquotedblleft{}roaring, boisterous\textquotedblright{}; CDIAL lists Punjabi \emph{dhun}, \textquotedblleft{}tune\textquotedblright{}, under it (CDIAL 6823).
\end{cousinweb}

\begin{grammarlens}[title={What you've built}]
One more word for this chapter.
\end{grammarlens}

\subsection*{Your turn}
\begin{itemize}
\raggedright
  \item \emph{Say it:} \emph{dhun}
  \item \emph{Say it:} \emph{dhun}, once more
  \item \emph{Say it:} \emph{us gīt dī dhun bahut miṭṭhī hai.} — The tune of that song is very sweet.
  \item \emph{Recall:} say the Punjabi for a musical instrument, then the Punjabi for a tabla, then say \emph{dhun} again
  \item \emph{Write it:} copy \textbf{\pa{ਧੁਨ}} once, checking the whole word
\end{itemize}

\begin{tcolorbox}[breakable,colback=teal!4,colframe=teal!35!black,title={Before you move on}]
What does \pa{ਧੁਨ} mean? (\textquotedblleft{}A tune; a melody\textquotedblright{}.) Say it once more.
\end{tcolorbox}
\section[\texorpdfstring{\pa{ਰਿਆਜ਼} (riāz)}{ਰਿਆਜ਼ (riāz)}]{\pa{ਰਿਆਜ਼} (riāz) — music practice; regular practice}

\label{lesson:PA-C348-riaz}



\begin{quote}
Before the new one: say the Punjabi for a tabla, then the Punjabi for a tune.
\end{quote}

\subsection*{You'll want to know: \pa{ਰਿਆਜ਼}}
\textbf{\pa{ਰਿਆਜ਼}} — \emph{riāz} — \textquotedblleft{}music practice; regular practice\textquotedblright{}.

\begin{quote}
\pa{ਮੇਰੀ} \pa{ਭੈਣ} \pa{ਸਵੇਰੇ} \pa{ਰਿਆਜ਼} \pa{ਕਰਦੀ} \pa{ਹੈ।}
\end{quote}

\begin{quote}
\emph{merī bhaiṇ savere riāz kardī hai.}
\end{quote}

\begin{quote}
My sister practises her music in the morning.
\end{quote}

A masculine noun.

\begin{cousinweb}
\textbf{\pa{ਰਿਆਜ਼}} (\emph{riāz}) is a loan, usually traced through Persian to Arabic \emph{riyāḍa}, \textquotedblleft{}exercise\textquotedblright{}.
\end{cousinweb}

\begin{grammarlens}[title={What you've built}]
One more word for this chapter.
\end{grammarlens}

\subsection*{Your turn}
\begin{itemize}
\raggedright
  \item \emph{Say it:} \emph{riāz}
  \item \emph{Say it:} \emph{riāz}, once more
  \item \emph{Say it:} \emph{merī bhaiṇ savere riāz kardī hai.} — My sister practises her music in the morning.
  \item \emph{Recall:} say the Punjabi for a tabla, then the Punjabi for a tune, then say \emph{riāz} again
  \item \emph{Write it:} copy \textbf{\pa{ਰਿਆਜ਼}} once, checking the whole word
\end{itemize}

\begin{tcolorbox}[breakable,colback=teal!4,colframe=teal!35!black,title={Before you move on}]
What does \pa{ਰਿਆਜ਼} mean? (\textquotedblleft{}Music practice; regular practice\textquotedblright{}.) Say it once more.
\end{tcolorbox}
\section[\texorpdfstring{\pa{ਗੂੰਜਣਾ} (gūnjṇā)}{ਗੂੰਜਣਾ (gūnjṇā)}]{\pa{ਗੂੰਜਣਾ} (gūnjṇā) — to echo; to ring out}

\label{lesson:PA-C348-gunjna}



\begin{quote}
Before the new one: say the Punjabi for a tune, then the Punjabi for music practice.
\end{quote}

\subsection*{You'll want to know: \pa{ਗੂੰਜਣਾ}}
\textbf{\pa{ਗੂੰਜਣਾ}} — \emph{gūnjṇā} — \textquotedblleft{}to echo; to ring out\textquotedblright{}.

\begin{quote}
\pa{ਤਬਲੇ} \pa{ਦੀ} \pa{ਆਵਾਜ਼} \pa{ਕਮਰੇ} \pa{ਵਿੱਚ} \pa{ਗੂੰਜਦੀ} \pa{ਹੈ।}
\end{quote}

\begin{quote}
\emph{table dī āvāz kamre vich gūnjdī hai.}
\end{quote}

\begin{quote}
The sound of the tabla rings out in the room.
\end{quote}

Past (it …): \pa{ਉਹ} \pa{ਗੀਤ} \pa{ਕਮਰੇ} \pa{ਵਿੱਚ} \pa{ਗੂੰਜਿਆ}, \emph{oh gīt kamre vich gūnjiā}.

\begin{cousinweb}
\textbf{\pa{ਗੂੰਜਣਾ}} (\emph{gūnjṇā}) is inherited from Sanskrit \emph{guñjati}, \textquotedblleft{}hums, buzzes\textquotedblright{}; CDIAL lists Punjabi \emph{gū̃jṇā}, \textquotedblleft{}to hum, buzz, growl\textquotedblright{}, under it (CDIAL 4175).
\end{cousinweb}

\begin{grammarlens}[title={What you've built}]
One more word for this chapter.
\end{grammarlens}

\subsection*{Your turn}
\begin{itemize}
\raggedright
  \item \emph{Say it:} \emph{gūnjṇā}
  \item \emph{Say it:} \emph{gūnjṇā}, once more
  \item \emph{Say it:} \emph{table dī āvāz kamre vich gūnjdī hai.} — The sound of the tabla rings out in the room.
  \item \emph{Recall:} say the Punjabi for a tune, then the Punjabi for music practice, then say \emph{gūnjṇā} again
  \item \emph{Write it:} copy \textbf{\pa{ਗੂੰਜਣਾ}} once, checking the whole word
\end{itemize}

\begin{tcolorbox}[breakable,colback=teal!4,colframe=teal!35!black,title={Before you move on}]
What does \pa{ਗੂੰਜਣਾ} mean? (\textquotedblleft{}To echo; to ring out\textquotedblright{}.) Say it once more.
\end{tcolorbox}
